\section{Perspectives d'amélioration et maturité DevSecOps}
\label{sec:perspectives-amelioration-maturite}

Cette section présente l'évaluation finale de la maturité DevSecOps de la plateforme \textit{SecureRAG Hub} et détaille la roadmap des améliorations à court, moyen et long terme. L'analyse critique s'appuie sur un audit technique rigoureux des configurations, des politiques d'admission et des mécanismes de détection système.

\subsection{Évaluation de la maturité et méthodologie d'audit}
\label{subsec:evaluation-maturite}

L'évaluation de la maturité DevSecOps a été conduite selon une méthodologie structurée couvrant cinq piliers fondamentaux : la gestion des identités et des secrets (IAM), la sécurité du code et de l'intégration continue (SAST/SCA/DAST), la sécurité de l'infrastructure et des conteneurs (IaC), le monitoring et la réponse aux incidents (SRE), ainsi que la documentation et la modélisation des menaces.

Cette approche multidimensionnelle attribue à l'environnement d'origine un score global de \textbf{78 sur 100}, le classant dans la catégorie « \textbf{Bon} » (environnement sécurisé mais perfectible). Les améliorations immédiates implémentées et configurées dans le cadre de cette itération visent à combler les principaux écarts identifiés, portant le score de maturité cible à \textbf{91 sur 100}, soit le niveau « \textbf{Excellent} » (défense en profondeur mature).

Le tableau~\ref{tab:score-maturite-devsecops} détaille la répartition de ce score et le gain estimé pour chaque pilier.

\begin{table}[htbp]
\centering
\small
\renewcommand{\arraystretch}{1.2}
\begin{tabular}{|l|c|c|c|l|}
\hline
\textbf{Pilier audit} & \textbf{Score initial} & \textbf{Gain} & \textbf{Score cible} & \textbf{Action clé associée} \\
\hline
IAM \& Secrets & 17 / 20 & +1 & 18 / 20 & Roadmap Vault + External Secrets \\
\hline
SAST / SCA / DAST & 18 / 20 & +2 & 20 / 20 & Intégration DAST OWASP ZAP \\
\hline
Supply Chain Security & 16 / 15 & — & 16 / 15 & SBOM et signature Cosign (bonus) \\
\hline
IaC \& Conteneurs & 14 / 20 & +4 & 18 / 20 & Kyverno Enforce + K8s Audit Policy \\
\hline
Monitoring \& Incidents & 8 / 15 & +6 & 14 / 15 & Falcosidekick + 5 alertes Prometheus \\
\hline
Documentation \& Menaces & 5 / 5 & — & 5 / 5 & Modèle STRIDE et runbooks complets \\
\hline
\textbf{TOTAL} & \textbf{78 / 100} & \textbf{+13} & \textbf{91 / 100} & \textbf{Niveau : Excellent} \\
\hline
\end{tabular}
\caption{Impact des améliorations sur la posture de sécurité globale}
\label{tab:score-maturite-devsecops}
\end{table}

\subsection{Axes d'amélioration immédiates (P0/P1 Implémentés)}
\label{subsec:axes-amelioration-immediates}

Les faiblesses identifiées lors de l'audit initial ont été traitées de manière proactive par des modifications directes dans le code et les descripteurs de déploiement de la plateforme :

\begin{enumerate}
  \item \textbf{Détection précoce des secrets (Shift-Left)} : L'ajout de la configuration \texttt{.pre-commit-config.yaml} permet d'exécuter localement \texttt{Gitleaks}, \texttt{ShellCheck} et des validateurs YAML avant chaque commit. Les secrets et erreurs de script sont ainsi bloqués sur la machine du développeur, avant même d'être poussés sur le dépôt distant.
  \item \textbf{Intégration du scan dynamique (DAST)} : Un stage \texttt{CD\_DAST} s'appuyant sur l'image officielle d'OWASP ZAP a été ajouté au pipeline de déploiement continu (\texttt{Jenkinsfile.cd}) et exposé via le \texttt{Makefile} (\texttt{make dast-baseline}). Le script de validation \texttt{validate-dast-report.sh} analyse le rapport JSON de ZAP et interrompt le pipeline (code retour non nul) en cas d'alerte de niveau critique ou élevé (\texttt{HIGH/CRITICAL}).
  \item \textbf{Routage actif des alertes Falco} : Le composant \texttt{Falcosidekick} a été activé dans \texttt{security/falco/values.yaml} afin d'extraire les alertes de détection d'intrusion au runtime (ex: reverse shell, exécution de binaire suspect) et de les acheminer en temps réel vers Loki et Alertmanager.
  \item \textbf{Enrichissement de l'alerting de sécurité} : Cinq nouvelles règles d'alerte Prometheus complexes ont été intégrées dans \texttt{prometheus-rules-security.yaml} :
  \begin{itemize}
    \item \texttt{ImagePullBackOff} : détection précoce d'une rupture de la \textit{supply chain} ou d'une signature révoquée.
    \item \texttt{ContainerHighCPU} : détection d'anomalies de consommation CPU pouvant indiquer un processus runaway ou du cryptomining.
    \item \texttt{SecretAccessAnomaly} : alerte en cas de pics suspects d'appels \texttt{GET/LIST} sur les secrets Kubernetes (tentative de récolte de credentials).
    \item \texttt{NamespaceQuotaNearLimit} : surveillance préventive de la saturation des ressources allouées.
    \item \texttt{DeploymentReplicasMismatch} : signalement de drift ou de blocage lors des déploiements applicatifs.
  \end{itemize}
  \item \textbf{Notification et receivers de production} : Les configurations stubs locales de l'Alertmanager ont été remplacées par une architecture de receivers réels (webhook + template Slack activable + template email) dans \texttt{alertmanager.yaml}.
\end{enumerate}

\subsection{Roadmap d'évolution et trajectoire cible (P2 Long Terme)}
\label{subsec:roadmap-evolution-p2}

Pour assurer la transition vers un environnement cloud de production hautement disponible et conforme aux standards industriels, une feuille de route en trois axes majeurs a été documentée :

\subsubsection{Gouvernance centralisée des secrets avec HashiCorp Vault}
La gestion actuelle basée sur SOPS+age, bien que robuste localement, est limitée pour une infrastructure multi-clusters en production. L'architecture cible documentée dans \texttt{vault-external-secrets-roadmap.md} préconise le déploiement d'un cluster \textbf{HashiCorp Vault} hautement disponible. Le couplage avec l'**External Secrets Operator (ESO)** permettra une synchronisation transparente et sécurisée des secrets depuis les moteurs de stockage clés-valeurs de Vault vers les objets \texttt{Secret} de Kubernetes, en s'appuyant sur des identités de comptes de service isolées.

\subsubsection{Authentification Supply Chain Keyless}
L'infrastructure Cosign actuelle repose sur une paire de clés cryptographiques statiques stockée dans Jenkins. La migration vers le mode \textbf{keyless} (sans clé statique), documentée dans \texttt{cosign-keyless-migration.md}, éliminera le risque de compromission de la clé privée en s'appuyant sur OpenID Connect (OIDC) et sur l'autorité de certification éphémère \textbf{Fulcio} de Sigstore. Les signatures d'images seront authentifiées à l'aide de l'identité OIDC du pipeline Jenkins ou GitHub Actions, archivée de manière transparente dans le registre public \textbf{Rekor}.

\subsubsection{Audit logging Kubernetes et Kyverno Enforce}
L'activation de l'audit logging natif de l'API Server Kubernetes, configurée dans \texttt{audit-policy.yaml}, permettra de tracer de manière exhaustive les appels sensibles (\texttt{exec}, \texttt{port-forward}, création de RBAC ou accès secrets) dans le namespace \texttt{securerag-hub}. Enfin, après une phase d'observation sans violation, les sept politiques Kyverno de base (actuellement en mode \texttt{Audit}) devront être basculées en mode \texttt{Enforce} via l'overlay \texttt{kyverno-enforce} préparé, interdisant le déploiement de tout manifest non conforme.

\subsection{Analyse SWOT de la posture DevSecOps}
\label{subsec:analyse-swot}

L'évaluation de la posture DevSecOps finale de \textit{SecureRAG Hub} peut être résumée à travers la matrice SWOT suivante :

\paragraph{Forces (Strengths).}
\begin{itemize}
  \item Socle de sécurité statique complet intégré au pipeline (Semgrep, Gitleaks, Trivy FS, Trivy Image).
  \item Supply chain robuste : signatures cryptographiques Cosign et génération systématique de SBOM CycloneDX.
  \item Cloisonnement réseau rigoureux (NetworkPolicy default-deny) et durcissement des conteneurs (Restricted Pod Security Standard).
  \item Preuves de déploiement et rapports d'audits générés et consolidés automatiquement.
\end{itemize}

\paragraph{Faiblesses (Weaknesses).}
\begin{itemize}
  \item Dépendance forte à un registre local d'images (\texttt{localhost:5001}) pour l'environnement local de démonstration (Kind).
  \item Risque d'usurpation en cas de compromission de la clé privée Cosign statique.
  \item Absence de persistance et de redondance géographique des bases de données applicatives dans le déploiement de base.
\end{itemize}

\paragraph{Opportunités (Opportunities).}
\begin{itemize}
  \item Migration keyless Sigstore / OIDC pour éliminer la gestion des clés cryptographiques.
  \item Automatisation complète de la rotation des secrets et de l'audit trail via HashiCorp Vault.
  \item Intégration d'un outil de prévention active runtime (IPS) basé sur eBPF (Tetragon) pour tuer automatiquement les processus malveillants.
\end{itemize}

\paragraph{Menaces (Threats).}
\begin{itemize}
  \item Vieillissement ou non-maintenance des règles de détection Falco ou Semgrep face à l'apparition de nouvelles CVE applicatives.
  \item Faux sentiment de sécurité si le mode d'admission Kyverno reste indéfiniment configuré en mode \texttt{Audit} plutôt qu'en mode \texttt{Enforce} en production.
\end{itemize}

\subsection{Conclusion de la démarche d'amélioration}
\label{subsec:conclusion-ameliorations}

La mise en œuvre des améliorations techniques (pre-commit, DAST, Falcosidekick, règles d'alertes avancées) et la formalisation des roadmaps architecturales démontrent que la sécurité de \textit{SecureRAG Hub} n'est pas traitée comme une contrainte statique, mais bien comme un processus itératif d'amélioration continue. Ce travail de consolidation referme les écarts opérationnels et offre au projet une posture DevSecOps robuste, industrialisée et directement alignée avec les exigences industrielles contemporaines de défense en profondeur.
